\documentclass[12pt]{article}
\usepackage[utf8]{inputenc}
\usepackage[T1]{fontenc}
\usepackage{hyperref}
\title{Sindihui Mixtec Language Resources for K-12 Curriculum Development}
\author{Compiled by Claude Sonnet 5.5 (claude-sonnet-5-5) for Andr\'as Kornai}
\date{}
\begin{document}
\maketitle

\section{Introduction}
Sindihui Mixtec (ISO 639-3 xts, Glottocode sind1277) is a nearly extinct Otomanguean language of Santa Maria (San Mateo) Sindihui, Oaxaca, Mexico. Wikipedia reports 34 native speakers in the 2005 census, all elderly. Joshua Project lists it as unwritten, and Global Recordings Network lists no recordings. Searches found no school programme, primer, dictionary or grammar.

\section{Primary Educational Resources}
None found.

\section{Linguistic and Reference Materials}
\begin{itemize}
\item Wikipedia entry: \url{https://en.wikipedia.org/wiki/Sindihui_Mixtec}
\item ASJP word-list entry (San Mateo Sindihui Mixtec): \url{https://asjp.clld.org/languages/SAN_MATEO_SINDIHUI_MIXTEC}
\item Wikipedia list of Oto-Manguean languages, for classification: \url{https://en.wikipedia.org/wiki/List_of_Oto-Manguean_languages}
\end{itemize}

\section{Religious and Community Resources}
Joshua Project language page: \url{https://joshuaproject.net/languages/xts}. Global Recordings Network page (no recordings shown): \url{https://globalrecordings.net/en/language/xts}.

\section{Digital and Audio Resources}
None found beyond the above.

\section{Curriculum Development Context}
Almost nothing is published. Any curriculum would first need community-led documentation: recording the few elderly speakers, deciding an orthography and building word lists, with the community's consent and leadership.
\end{document}